\chapter{Zakończenie}

\todo[inline]{Sugerowałbym zmianę nazwy rozdziału na Zakończenie, co jest to zgodne z wymaganiami prac magisterskich (proszę dodatkowo przejrzeć czy nie pomnieliśmy czegoś bezpośrednio z tabelki w kontekście całej pracy - link w komentarzu)} %https://wit.pwr.edu.pl/studenci/dyplomanci/praca-dyplomowa

Niniejsza praca koncentrowała się na analizie problemu uczenia głębokich sieci neuronowych w warunkach ograniczonej liczby danych treningowych – skutecznego uczenia głębokich sieci neuronowych w reżimie małego zbioru danych. Na przykładzie trudnego zadania detekcji gęsto upakowanych obiektów (worków z piaskiem), zaprojektowano i wdrożono zautomatyzowane środowisko eksperymentalne umożliwiające optymalizację modeli. Przeprowadzone badania pozwoliły na obiektywne porównanie wpływu metod \textit{data-centric} – iteracyjnego uczenia pół-nadzorowanego (SSL) oraz syntetycznej augmentacji danych – na nowoczesne architektury detekcyjne: splotowe (YOLOv8n, YOLOv26n) oraz oparte na mechanizmach uwagi (RT-DETR-L).

\section{Wnioski z przeprowadzonych badań}

Eksperymenty przeprowadzone w ramach pracy umożliwiły sformułowanie szeregu istotnych wniosków, odpowiadających na postawione cele badawcze. Główne obserwacje z podziałem na zastosowane metody prezentują się następująco:

\begin{itemize}
    \item \textbf{Ograniczenia modeli bazowych i zjawisko przeuczenia:} Trenowanie nowoczesnych architektur wyłącznie na małym zbiorze danych (ok. 210 próbek) może prowadzić do problemów z generalizacją. Modele splotowe (YOLO) wykazały tendencję do generowania fałszywie pozytywnych detekcji (halucynacji w obszarach tła). Z kolei architektura oparta na transformerach (RT-DETR-L) wykazała dużą wrażliwość na ograniczoną liczbę danych treningowych – model ten wykazywał oznaki szybkiego przeuczenia już w początkowej fazie treningu, a jego krzywa zbieżności była wysoce niestabilna.
    
    \item \textbf{Niejednoznaczność pseudo-etykietowania (Metoda A):} Iteracyjne uczenie pół-nadzorowane z wykorzystaniem naturalnych obrazów nieoznakowanych nie jest rozwiązaniem uniwersalnym. Dla wysoce zoptymalizowanego modelu YOLOv8n włączenie pseudo-etykiet w proces optymalizacji przyniosło spadek precyzji, co było efektem błędu potwierdzenia (\textit{confirmation bias}). Z drugiej strony, dla modelu YOLOv26n metoda ta pełniła funkcję dodatkowego mechanizmu regularyzacji, poprawiając globalną metrykę mAP. Dodatkowo, dostarczenie dodatkowych danych znacząco ustabilizowało proces treningu modelu RT-DETR-L, wyraźnie opóźniając zjawisko jego przeuczenia.
    
    \item \textbf{Skrajne reakcje na dane syntetyczne (Metoda B):} Wykorzystanie modeli dyfuzyjnych do generowania syntetycznych próbek negatywnych (czystych teł) uwidoczniło różnice pomiędzy analizowanymi architekturami. Wymuszenie na modelu YOLOv26n uniezależnienia się od kontekstu otoczenia przyniosło korzystne rezultaty – zredukowano liczbę halucynacji, opóźniono przeuczenie o ponad 200 epok i osiągnięto najwyższy wzrost wydajności w całych badaniach (+3.4 p.p. mAP@0.5). 
    
    \item \textbf{Krytyczna wrażliwość transformerów wizyjnych:} Jednym z bardziej interesujących wyników badań jest zachowanie modelu RT-DETR-L wobec danych syntetycznych. Wprowadzenie wygenerowanych teł doprowadziło do istotnego pogorszenia wyników uzyskiwanych przez ten model. Wyniki sugerują, że sieci oparte na mechanizmach uwagi, pozbawione uprzedzeń indukcyjnych (\textit{inductive biases}) znanych z klasycznych sieci CNN, mogą być bardziej wrażliwe na przesunięcie domeny i nienaturalne artefakty wprowadzane przez modele generatywne.
\end{itemize}

Osiągnięte wyniki wskazują, że w dziedzinie głębokiego uczenia nie istnieje uniwersalne rozwiązanie problemu niedoboru danych. Skuteczność nowoczesnych technik augmentacji i uczenia pół-nadzorowanego w dużym stopniu zależy od specyfiki architektury wybieranego modelu oraz charakterystyki jego bazowych problemów. 

Uzyskane wyniki pozwalają na częściowe potwierdzenie postawionej hipotezy badawczej. Zastosowanie dodatkowych danych, zarówno rzeczywistych w procesie pseudo-etykietowania, jak i syntetycznych w ramach augmentacji, przyczyniło się do poprawy zdolności generalizacji wybranych modeli. Jednocześnie wykazano, że skuteczność tych metod jest silnie uzależniona od architektury modelu, co potwierdza\todo{bardziej wskazuje/sugeruje -- tutaj proszę, aby sprawdził Pan też w innych miejscach. Należy być ostrożnym, aby nie twierdzić (szczególnie z ogromną stanowczością) czegoś, na co nie ma Pan podparcia w swoich danych/eksperymentach. Bezpieczniej i łagodniej, zawsze może Pan podkreślać, że zakres pracy był ograniczony - np. do danego zbioru danych, więc nie możemy wyciągać zbyt ogólnych wniosków; ale też chwalić swoje mocne strony: że przeprowadzenie eksperymentów na wielu modelach, choć utrudniło przeprowadzenie badań, ale pozwala na wyciągniecie ciekawych wniosków...}, że nie istnieje jedno uniwersalne rozwiązanie skuteczne dla wszystkich analizowanych detektorów.

\section{Możliwe kierunki dalszego rozwoju}

Wyniki uzyskane w niniejszej pracy stanowią solidny fundament do dalszych eksploracji w obszarze detekcji obiektów w reżimie małych zbiorów danych. Obserwacje zebrane podczas analizy eksperymentów pozwalają na wyznaczenie obiecujących kierunków przyszłych badań:

\begin{enumerate}
    \item \textbf{Skalowanie objętości danych nieoznakowanych i syntetycznych:} W przeprowadzonych badaniach skutecznie zweryfikowano mechanikę potoku SSL na kontrolowanej, stosunkowo niewielkiej puli danych dodatkowych (60 kafelków tła i 60 obrazów bez etykiet). Naturalnym kierunkiem dalszych badań jest drastyczne zwiększenie rozmiaru tych zbiorów. Wykorzystanie np. tysięcy surowych zdjęć lub masowo wygenerowanych teł syntetycznych mogłoby zwiększyć skuteczność asymetrycznej augmentacji w iteracyjnych pętlach sprzężenia zwrotnego, skutkując znaczną poprawą zdolności do generalizacji.
    \item \textbf{Zaawansowane filtrowanie pseudo-etykiet:} Błąd potwierdzenia zauważalny u części modeli można zniwelować, zastępując sztywne progowanie pewności (wykorzystane w tej pracy) metodami estymacji niepewności modelu (ang. \textit{uncertainty estimation}), takimi jak \textit{Monte Carlo Dropout} czy wykorzystanie zespołu modeli (\textit{ensembling}).
    \item \textbf{Rozwój potoku generatywnego:} W obliczu niestabilności modelu RT-DETR na sztucznie wygenerowanych tłach, wartościowym kierunkiem byłoby zastosowanie technik adaptacji domeny (ang. \textit{domain adaptation}) lub wykorzystanie nowszych, nowszych modeli dyfuzyjnych (np. Stable Diffusion XL, Flux) do generowania próbek pozbawionych mikroskopijnych artefaktów wizyjnych. Alternatywą jest zastosowanie technik \textit{inpaintingu}, polegających na wklejaniu syntetycznych obiektów w naturalne tła, zamiast generowania całych teł.
    \item \textbf{Dynamiczne wagowanie funkcji straty w SSL:} Obecne środowisko można rozbudować o algorytmy dynamicznie dostosowujące wpływ komponentu straty ze zbioru nieoznakowanego (tzw. \textit{unsupervised loss weight}) w zależności od epoki treningu. Pozwoliłoby to na łagodniejsze wprowadzanie pseudo-etykiet w początkowych, wysoce niestabilnych fazach nauki modeli opartych na transformerach.
    \item \textbf{Ewaluacja w środowiskach wieloklasowych:} Zweryfikowanie poprawności wyciągniętych wniosków (szczególnie w kontekście podatności architektur) na bardziej złożonych, wieloklasowych zestawach danych oraz na obrazach pochodzących z innych domen, np. z kamer termowizyjnych lub obrazowania medycznego.
\end{enumerate}

W ramach pracy opracowano zautomatyzowane środowisko eksperymentalne oraz przeprowadzono analizę wpływu badanych metod, które mogą bezpośrednio posłużyć inżynierom przy projektowaniu lepszych \todo{za silne w tym kontekście, wystarczy: 'lepszych'} systemów wizyjnych w reżimie małego zbioru danych.\todo{bliżej pracy bardziej w reżimie małego zbioru danych, dodatkowo pozwoliłbym tutaj na }.